%****************************** Preambulo ******************************%
\documentclass[letterpaper, 11pt]{article} 

\usepackage[spanish, es-tabla, es-nodecimaldot]{babel}
	% spanish, cambia el idioma del documento de ingles a español. 
	% es-tabla, traduce "table" como tabla en lugar de cuadro (cuadro se utiliza en castellano).
	% es-nodecimaldot, cambia la coma decimal por punto decimal en los entornos matemáticos. Ejemplo: cambia 1,23 por 1.23.

\usepackage[utf8]{inputenc}		% Paquete que soluciona los problemas de acentos y la inserción de la ñ.
\usepackage{graphicx}			% Paquete necesario para insertar imágenes.
\usepackage{float}				% Paquete que corrige la ubicación de imágenes. \begin{figure}[H]

\usepackage{appendix}
\renewcommand{\appendixname}{Apéndices}
\renewcommand{\appendixtocname}{Apéndices}
\renewcommand{\appendixpagename}{Apéndices}


\usepackage[vmargin=2.5cm,hmargin=2.5cm]{geometry} % Asigna tamaños a los margenes de todo el documento.
	% vmargin (vertical margin), asigna un tamaño de margen superior e inferior del mismo tamaño.
	% hmargin (horizontal margin), asigna un tamaño de margen al lado izquierdo y derecho del mismo tamaño.
	
\usepackage[small]{caption}		% Reduce el tamaño de letra utilizado en los pies de figura.
\usepackage{hyperref}		% Agrega enlaces internos de las secciones, figuras y tablas.
\hypersetup{colorlinks=true, linkcolor=[rgb]{0.2,0.2,0.2}, citecolor=[rgb]{0.2,0.2,0.2}}
% colorlinks - permite modificar el color de los enlaces.
% linkcolor - se asigna el color deseado a los enlaces de figuras y tablas. 
% citecolor - se asigna el color deseado a las citas y referencias.

%******************************* Portada *******************************%
\title{		% Comienza entorno título
\centering
\includegraphics[scale=.8]{imagenes/UDG_byn.png}\\
\Large UNIVERSIDAD DE GUADALAJARA\\
\small CENTRO UNIVERSITARIO DE CIENCIAS EXACTAS E INGENIERÍAS\\
\vspace{20pt}
\large SISTEMAS DE MEDICIÓN\\[1cm]
\textbf{\huge P\# - Nombre de la práctica}\\[0.5cm]}	% Termina entorno título
% P\# Es el número de la práctica P1, P2, P3, ...
% El nombre de la práctica se encuentra en las instrucciones de cada práctica.
% Por ejemplo: P2 - Incertidumbre de elementos eléctricos  

%******************************* Autores *******************************%
\author{								
\centering
\begin{tabular}{ll}
Alumnos:			&	Nombre de alumno 1, (Código)\\
					&	Nombre de alumno 2, (Código)\\
Profesor:			&	Dr. Guillermo García Torales
\end{tabular}
}
\date{\today}	% Se genera la fecha del día en que se compila

%****************************** Documento ******************************%
\begin{document}

\pagestyle{empty}
\maketitle	% Genera la portada con base a lo definido en \title, \author y \date.
\thispagestyle{empty}
\newpage	% Salto de página.

\pagestyle{plain}
\pagenumbering{Roman}	% Cambia la númeración de la página a números romanos.
\tableofcontents	% Genera una tabla de contenido o índice.
\newpage

\section*{Lista de símbolos} % Cambia la tipografía de las letras que están dentro del llaves,
% como si fuera un título dentro del documento pero sin numerarlo.
\begin{tabular}{p{5mm} c p{120mm} }
\multicolumn{3}{l}{Mayúsculas}\\
\\
$T$ & --- & Temperatura (ºC).\\
$V$ & --- & Volumen (m$^3$).\\
$K$ & --- & Constante de Boltzman (J/ºK)
\\
\\
\multicolumn{3}{l}{Minúsculas}\\
\\
$c$ & --- & Velocidad de la luz (m/s). La velocidad de la radiación electromagnética es independiente de la velocidad del emisor.\\
$g$	& ---	& Aceleración gravitacional (m/s$^2$).\\
$i$ & --- & Raíz de menos uno.\\
\\
\\
\multicolumn{3}{l}{Letras griegas}\\
\\
$\tau$ & --- & Constante de tiempo (s).\\
$\alpha$ & --- & Constante de amortiguamiento.\\
$\omega$ & --- & Frecuencia angular (rad/s).
\end{tabular}
\newpage
\pagenumbering{arabic}	% Cambia la númeración de la página a números arábigos.

\begin{abstract}
El principal propósito del resumen es proporcionar al lector una descripción del proyecto y sus resultados. Debe ser breve (150 palabras aproximadamente) e informativo. Este resumen, debe proporcionar la información necesaria para que el lector decida si es, o no, de su interés. A partir de este punto, es importante que considere, que un reporte técnico debe escribirse de manera impersonal y en pasado. De tal manera, que debemos evitar palabras como ``calculé'', ``medi'', ``observe''; siendo las palabras correctas ``se calculó'', ``se midió'', ``se  observo''. 
\end{abstract}

\section{Introducción}
La introducción tiene como propósito, poner en contexto el tema principal sobre el que se basa el experimento o práctica a realizar. Por lo que se debe dar al lector una breve secuencia descriptiva que va de lo general a lo específico. Se mencionan los antecedentes y relevancia de este experimento con base en sus aplicaciones. Se hace referencia a trabajos previamente publicados relevantes para el experimento y se muestra como es que se aplica a nuestro  experimento en particular.\\
Además, se puede agregar un breve párrafo de transición antes de abordar el marco teórico así como proporcionar una guía para el lector sobre lo que va a seguir en el informe.

\section{Marco teórico}
Se deben exponer los fundamentos teóricos principales para realizar el experimento. Ésta base teórica sirve para dar conformidad a los resultados siguiendo un método científico, sustentando los resultados con las teorías expuestas.\\
Es importante mencionar y describir la expresión matemática del modelo o ley que se va a utilizar, explicando de manera clara y breve cada uno de sus  términos.\\
La teoría debe ser explicada en su  totalidad pero de manera breve, pasando por alto teorías fundamentales que ya son conocidas o que no son tan relevantes  para la experimentación.

\section{Desarrollo}
En esta sección se debe describir brevemente los métodos experimentales, analíticos y numéricos utilizados para llegar a los resultados. Debe realizarse de forma detallada, suficiente para permitir una evaluación crítica de los métodos utilizados y posibles réplicas de los resultados obtenidos realizados por un tercero. Debes realizar estimaciones de la incertidumbre para todos los parámetros y los procedimientos utilizados para llegar a los resultados que se presenten. Un diagrama de bloques o esquema, pueden ser de gran utilidad para dar un mejor enfoque al lector.  Un diagrama de flujo ayuda a describir un programa realizado en algún dispositivo, pero es importante considerar que el código de programación no debe estar en el informe, en caso de querer agregar el código  se puede anexar  en la parte final del documento como un apéndice.

\section{Resultados}
Los resultados de tu experimento deben ser presentados en esta sección, usualmente se utilizan gráficos, figuras y tablas para facilitar su entendimiento. Los datos encontrados en el experimento son presentados pero nunca discutidos o evaluados en esta sección. Observamos la Tabla \ref{tab:tabla1} Mantén en mente que lo que se desea es presentar los resultados de una manera concisa. Pon tus resultados en una gráfica siempre que esto sea posible. Debes mencionar de forma específica el contenido en todas las figuras y tablas. Nunca se debe de decir “por ejemplo”, es mejor decir “los resultados obtenidos se muestran en la figura X”.
\begin{table}
\centering
\caption{Tabla de ejemplo.}
\label{tab:tabla1}
\begin{tabular}{|l|c|c|}
\hline
Palabra E$_{1,1}$	&	E$_{1,2}$	&	E$_{1,3}$ \\ \hline
E$_{2,1}$	&	E$_{2,2}$	&	E$_{2,3}$ \\ \hline
E$_{3,1}$	&	E$_{3,2}$	&	E$_{3,3}$ \\ \hline
\end{tabular}
\end{table}

\section{Conclusiones}
Escribir un pequeño resumen de los principales hallazgos del experimento. No introducir nada nuevo ni continuar con la discusión de los resultados. Centrar los comentarios en cuestiones técnicas y no de carácter personal.

\section{Referencias}
Libros \\
1 Author Initial. Author Surname, Title, Edition, (Publisher, City, Year Published), p. Pages Used.
Ejemplo: R.A. Serway and J.W. Jewett, Physics for Scientists and Engineers, 6th Ed. (Thomson, Belmont, CA, 2004), pp. 100-102.


Revistas \\
2 Author Initial. Author Surname, journal, Publication Title Volume number, (Year Published).
Ejemplo: J.J. Amato et al., Phys. Rev. Lett. 21, 1709 (1968). 

Documento de Internet\\
3 Author Initial. Author Surname, Title, WWW Document, (Internet URL).\\
Ejemplo: T.M. Huber, How To Locate Material for Formal Reports, WWW Document, (http://physics.gac.edu/~huber/misc/finding.htm)



\appendix
\clearpage
\addappheadtotoc
\appendixpage

\section{Notas y firma}
Anexen los escaneos o en su defecto, fotografías claras y de buena resolución de las anotaciones que hicieron durante la práctica junto con la práctica.\\
Sin firma no se revisa la práctica.

\section{Otros}
En el apéndice deben estar todas aquellas herramientas o elementos adicionales que no  son relevantes para el informe pero si para realizar la práctica o experimento. La finalidad de  un reporte es que el lector, si así lo desea, pueda reproducir el experimento al cien por ciento. Por ejemplo se pueden agregar:
\begin{itemize}	% Comienza el entorno para agregar viñetas, cada \item corresponde a una nueva viñeta.
\item Hojas de datos.
\item Código.
\item Fotografías descriptivas del arreglo experimental.
\item Especificaciones del equipo o  material utilizado.
\item Información completa acerca del software utilizado.
\end{itemize}
Cada uno de estos debe estar en un apéndice  distinto, como aparece en este documento, por ejemplo supongamos que es necesario agregar información de hojas de datos  y código, tendría tres apéndices:\\[4mm]
\large Apéndices
\normalsize
\begin{itemize}
\item[A.] Notas y firma
\item[B.] Hoja de datos del diodo 1N4007
\item[C.] Código de programación de Arduino de la conversión de voltaje a presión
\end{itemize}
\noindent
\\[3mm]
Los reportes no serán revisados en los siguientes casos:
\begin{itemize}
\item Que el nombre del documento no corresponda  al establecido en clase. \\SM-Sección-Apellidos-P\#.
\item Que el reporte no contenga el escaneo de las hojas donde se encuentra la firma.
\item Que el reporte no esté realizado en \LaTeX .
\item Que el reporte no este con el formato establecido.
\item Que el reporte se suba posterior a la fecha de entrega.
\item Que el reporte no  contenga nombres de los  alumnos.
\end{itemize}
% Después de cerrar el entorno documento con \end{document} se finaliza todo el documento de texto
% creado, por lo que no  puede haber nada después de este comando ya que será marcado como error.
\end{document}
